%% kapa: exam=1 round=29 year=2018 period=1 form=A subject=economics subject_ko=경제학원론 seq=13 number=53
%% source: Q-Net 감정평가사 기출문제 (cst00309) artlSeq=5202808 file=extracted/1008/5202808/2018년 제29회 감정평가사 1차시험 기출문제/2018년 제29회 감정평가사 1차시험 기출문제/1교시 A형.pdf
%% conversion: kordoc 4.24.51 -> tools/gen1.py (auto); review=proofread-sonnet; flags=tikz
\begin{question}{53}
소비자 甲이 두 재화 X, Y를 소비하고 효용함수는 $U(x,y)=\min\{x+2y,2x+y\}$ 이다. 소비점 (3,3)을 지나는 무차별곡선의 형태는? (단, $x$는 X의 소비량, $y$는 Y의 소비량이다.)
\begin{center}
\begin{tabular}{@{}l@{\hspace{1.5em}}l@{\hspace{1.5em}}l@{}}
\begin{tikzpicture}[x=0.1mm,y=-0.1mm,font=\scriptsize,line width=0.7pt]
\node[anchor=north east,font=\small] at (20,10) {①};
\draw[-{Stealth[length=1.4mm]}] (36,305) -- (36,10);
\draw[-{Stealth[length=1.4mm]}] (36,305) -- (335,305);
\node at (22,6) {$\boldsymbol{y}$};
\node at (326,320) {$\boldsymbol{x}$};
\node at (20,322) {0};
\node at (14,163) {3};
\node at (180,322) {3};
\draw[dashed] (36,163) -- (180,163) -- (180,305);
\draw (52,37) -- (308,291);
\end{tikzpicture} &
\begin{tikzpicture}[x=0.1mm,y=-0.1mm,font=\scriptsize,line width=0.7pt]
\node[anchor=north east,font=\small] at (20,10) {②};
\draw[-{Stealth[length=1.4mm]}] (36,305) -- (36,10);
\draw[-{Stealth[length=1.4mm]}] (36,305) -- (335,305);
\node at (22,6) {$\boldsymbol{y}$};
\node at (326,320) {$\boldsymbol{x}$};
\node at (20,322) {0};
\node at (14,163) {3};
\node at (180,322) {3};
\draw[dashed] (36,163) -- (180,163) -- (180,305);
\draw[domain=1.5:5.9,samples=40,variable=\t] plot ({36+48*\t},{305-47*(9/\t)});
\end{tikzpicture} &
\begin{tikzpicture}[x=0.1mm,y=-0.1mm,font=\scriptsize,line width=0.7pt]
\node[anchor=north east,font=\small] at (20,10) {③};
\draw[-{Stealth[length=1.4mm]}] (36,305) -- (36,10);
\draw[-{Stealth[length=1.4mm]}] (36,305) -- (335,305);
\node at (22,6) {$\boldsymbol{y}$};
\node at (326,320) {$\boldsymbol{x}$};
\node at (20,322) {0};
\node at (14,163) {3};
\node at (180,322) {3};
\draw[dashed] (36,163) -- (180,163) -- (180,305);
\draw (180,31) -- (180,163) -- (312,163);
\end{tikzpicture}\\
\begin{tikzpicture}[x=0.1mm,y=-0.1mm,font=\scriptsize,line width=0.7pt]
\node[anchor=north east,font=\small] at (20,10) {④};
\draw[-{Stealth[length=1.4mm]}] (36,305) -- (36,10);
\draw[-{Stealth[length=1.4mm]}] (36,305) -- (335,305);
\node at (22,6) {$\boldsymbol{y}$};
\node at (326,320) {$\boldsymbol{x}$};
\node at (20,322) {0};
\node at (14,163) {3};
\node at (180,322) {3};
\draw[dashed] (36,163) -- (180,163) -- (180,305);
\draw (110,29) -- (180,163) -- (314,230);
\end{tikzpicture} &
\begin{tikzpicture}[x=0.1mm,y=-0.1mm,font=\scriptsize,line width=0.7pt]
\node[anchor=north east,font=\small] at (20,10) {⑤};
\draw[-{Stealth[length=1.4mm]}] (36,305) -- (36,10);
\draw[-{Stealth[length=1.4mm]}] (36,305) -- (335,305);
\node at (22,6) {$\boldsymbol{y}$};
\node at (326,320) {$\boldsymbol{x}$};
\node at (20,322) {0};
\node at (14,163) {3};
\node at (180,322) {3};
\draw[dashed] (36,163) -- (180,163) -- (180,305);
\draw (48,99) -- (180,163) -- (245,290);
\end{tikzpicture}\\
\end{tabular}
\end{center}

\end{question}
\answer{4}
